\documentclass[11pt]{article}
\usepackage[margin=1in]{geometry}
\usepackage{amsmath,amssymb,mathtools}
\usepackage{fontspec,unicode-math}
\setmainfont{Latin Modern Roman}
\setsansfont{Latin Modern Sans}
\setmathfont{Latin Modern Math}
\usepackage{xurl,hyperref}
\hypersetup{colorlinks=true,urlcolor=blue,linkcolor=blue}
\newcommand{\sref}[2]{\hyperlink{source:#1:#2}{[S#2]}}
\newcommand{\cataloguescope}[1]{\paragraph{Recorded mathematical question.}#1}
\begin{document}
\section{Crossings in random square tilings}\label{statement}

\subsection{Definitions and mathematical statement}

\paragraph{Shared notation.}$\mathbb N=\{1,2,\ldots\}$, and $\mathbb Z,\mathbb Q,\mathbb R,\mathbb C$ denote the integers, rationals, reals and complex numbers. Any other conventions are specified locally.


Let $\mathcal T$ be a finite or countable tiling of the closed unit square by squares of positive side length, with disjoint interiors and with at most three squares meeting at any corner. Color its squares black independently with probability $1/2$. A black left-to-right crossing is a connected black chain, or connected union of black tiles in the countable case, meeting both vertical sides. Write $q(\mathcal T)$ for its probability and $m(\mathcal T)=\sup_{Q\in\mathcal T}\operatorname{diam}(Q)$ for the mesh.
\paragraph{Statement recorded in the supplied source.}
Prove that there is an absolute $c>0$ such that $q(\mathcal T)>c$ for every such tiling. Moreover, for every sequence of such tilings $\mathcal T_n$ with $m(\mathcal T_n)\to0$, prove $q(\mathcal T_n)\to1/2$. \sref{statement}{2}
\subsection{Short English statement}
Show that random square tilings always have a uniformly positive horizontal crossing probability and approach one half as the largest tile shrinks.
\subsection{Sources}
\begin{description}
\item[\hypertarget{source:statement:1}{S1}] ULAM.AI and the UnsolvedMath Contributors, \emph{UnsolvedMath: A Curated Collection of Open Mathematics Problems}, catalogue entry. CC BY 4.0. \url{https://huggingface.co/datasets/ulamai/UnsolvedMath}.
\item[\hypertarget{source:statement:2}{S2}] Itai Benjamini, \emph{Uniformization and percolation}, October 2015; original author slides preserved in Arquivo.pt capture of 31 December 2020, PDF slides 22--23, Motivation for the conjecture, a short detour. \url{https://arquivo.pt/noFrame/replay/20201231041533id_/http://www.wisdom.weizmann.ac.il/~itai/conformaltalk1.pdf}.
\end{description}
\label{end:statement}
\end{document}
